\section{Cel i zakres dokumentu}

Niniejszy dokument stanowi kompletną dokumentację techniczną systemu \textit{ForestFire} -- rozproszonej platformy do symulacji pożarów lasu oraz wspomagania decyzji o działaniach gaśniczych i patrolowych. Zebrano w nim w jednym miejscu opis architektury, logiki domenowej i algorytmów decyzyjnych, kontraktów komunikacyjnych między modułami oraz obsługi systemu od strony użytkownika.

Dokument obejmuje:
\begin{itemize}
    \item architekturę systemu, jego komponenty oraz przepływ danych (rozdział~2),
    \item model symulacji pożaru oparty na parametrycznym L-systemie, detekcję, algorytm MCTS oraz heurystyki decyzyjne (rozdział~3),
    \item interfejsy REST, kontrakty RabbitMQ i formaty wymiany danych JSON (rozdział~4),
    \item obsługę systemu, interfejs wizualizacji oraz znane ograniczenia (rozdział~5).
\end{itemize}

\subsection{Przeznaczenie}

Dokument adresowany jest do zespołu rozwojowego i utrzymaniowego, przyszłych kontynuatorów projektu, integratorów systemów oraz audytorów technicznych. Zakłada się u czytelnika podstawową znajomość systemów rozproszonych i kolejek wiadomości.
